% GENERATED FILE. Edit canonical lessons, then run npm run generate:books.
% canonical-source-hash: fnv1a64:8eddbc847cfe04f9
% canonical-lessons: ZH-C39-he, ZH-W39-he-drink, ZH-C39-shui, ZH-W39-shui, ZH-C39-ke, ZH-C39-niunai, ZH-C39-cha

\chapter{Drinking}
\label{ch:zh-drinking}
\hlchaptermodality{39}

\begin{chapteropening}
\textbf{By the end of this chapter:} \emph{I can say what I drink and ask somebody what they would like.}

Complete the last lesson of chapter 39: I can say what I drink and ask somebody what they would like.
\end{chapteropening}

\section[\texorpdfstring{\zh{喝} (hē)}{喝 (hē)}]{\zh{喝} (hē) — to drink}

\label{lesson:ZH-C39-he}



\begin{quote}
Before the new one: say the Mandarin for hungry, then the Mandarin for dish of food.
\end{quote}

\subsection*{You'll want to know: \zh{喝}}
\textbf{\zh{喝}} — \emph{hē} — \textquotedblleft{}to drink\textquotedblright{}.

\begin{quote}
\textbf{\zh{你喝什么}} \emph{nǐ hē shénme} — What would you like to drink?
\end{quote}

The new character, in pieces: \zh{喝} — \zh{口} on the left; on the right \zh{日} on top of a wrap with a small shape inside it. The next lesson puts it in your hand.

\begin{cousinweb}
\textbf{\zh{喝}} (\emph{hē}) is an old character: the Guangyun records it as \emph{xat} (Guangyun, in Baxter's transcription), glossed 'to shout at'. The sense 'to drink' is not given there, so how it came to mean drink is not settled.
\end{cousinweb}

\begin{grammarlens}[title={What you've built}]
One more word for everyday answers.
\end{grammarlens}

\subsection*{Your turn}
\begin{itemize}
\raggedright
  \item \emph{Say it:} \emph{hē}
  \item \emph{Say it:} \emph{hē}, once more
  \item \emph{Say it:} \emph{nǐ hē shénme} — What would you like to drink?
  \item \emph{Recall:} say the Mandarin for hungry, then the Mandarin for dish of food, then say \emph{hē} again
\end{itemize}

\begin{tcolorbox}[breakable,colback=teal!4,colframe=teal!35!black,title={Before you move on}]
What does \zh{喝} mean? (\textquotedblleft{}To drink\textquotedblright{}.) Say it once more.
\end{tcolorbox}
\section[\texorpdfstring{\zh{喝} (hē)}{喝 (hē)}]{\zh{喝} — a mouth, a sun and a wrap}

\label{lesson:ZH-W39-he-drink}



\begin{quote}
The word first, then the character inside it.

\begin{itemize}
\raggedright
  \item \emph{Recall:} say \emph{hē} — \textbf{R1}, one lesson back, before it decays
\end{itemize}

\begin{quote}
\zh{口} \zh{日} \zh{吃}
\end{quote}

\emph{Read it:} each one aloud, with nothing said first — \emph{kǒu}, \emph{rì}, \emph{chī}
\end{quote}

\subsection*{Script — a mouth, a sun and a wrap}
\hlblockfigure{\detokenize{figures/ZH-W39-he-drink-filmstrip.pdf}}{How \zh{喝} is written, stroke by stroke}

\begin{quote}
\zh{喝}
\end{quote}

\textbf{Three pieces: \zh{口} on the left, then \zh{日}, then the wrap under it.} Do not count to twelve; count three.

\textbf{\zh{口}, three strokes}, small, as in \zh{吃}.

\textbf{\zh{日}, four strokes}, short and wide, at the top right.

\textbf{The wrap, five strokes.} A short left-falling stroke. A horizontal that turns down and hooks, making a roof. Inside it, a small left-falling stroke and a dot. Then a corner stroke, down and along, closing them in.

3 + 4 + 5 = twelve, and the first seven are two shapes you know.

\subsection*{Your turn}
\begin{itemize}
\raggedright
  \item \emph{Write it:} \zh{喝}, all twelve, with the model in front of you
  \item \emph{Cover:} write it again from memory, and count the pieces, not the strokes
\end{itemize}

\begin{tcolorbox}[breakable,colback=teal!4,colframe=teal!35!black,title={Before you move on}]
How many pieces in \zh{喝}? (\textbf{Three} — \zh{口}, \zh{日}, and the wrap.) How many strokes before the wrap starts? (\textbf{Seven}.) What goes inside the wrap? (\textbf{A small left-falling stroke and a dot}.)

Source: \href{https://raw.githubusercontent.com/chanind/hanzi-writer-data/68d10a4b21150cae5e1ebbd223eed289cf32d90c/data/\%E5\%96\%9D.json}{Hanzi Writer Data \zh{喝}, pinned stroke paths and medians}.
\end{tcolorbox}
\section[\texorpdfstring{\zh{水} (shuǐ)}{水 (shuǐ)}]{\zh{水} (shuǐ) — water}

\label{lesson:ZH-C39-shui}



\begin{quote}
Before the new one: say the Mandarin for dish of food, then the Mandarin for to drink.
\end{quote}

\subsection*{You'll want to know: \zh{水}}
\textbf{\zh{水}} — \emph{shuǐ} — \textquotedblleft{}water\textquotedblright{}.

\begin{quote}
\textbf{\zh{我喜欢喝水}} \emph{wǒ xǐhuan hē shuǐ} — I like drinking water.
\end{quote}

The new character, in pieces: \zh{水} — a hooked vertical down the middle, one stroke on its left that runs across and falls away, and two strokes on its right. The next lesson puts it in your hand.

\begin{cousinweb}
\textbf{\zh{水}} (\emph{shuǐ}) is inherited from Middle Chinese \emph{sywijX} (Guangyun, in Baxter's transcription), which the dictionary explains as one of the five elements, the one of the north, and as 'level'.
\end{cousinweb}

\begin{grammarlens}[title={What you've built}]
One more word for everyday answers.
\end{grammarlens}

\subsection*{Your turn}
\begin{itemize}
\raggedright
  \item \emph{Say it:} \emph{shuǐ}
  \item \emph{Say it:} \emph{shuǐ}, once more
  \item \emph{Say it:} \emph{wǒ xǐhuan hē shuǐ} — I like drinking water.
  \item \emph{Recall:} say the Mandarin for dish of food, then the Mandarin for to drink, then say \emph{shuǐ} again
\end{itemize}

\begin{tcolorbox}[breakable,colback=teal!4,colframe=teal!35!black,title={Before you move on}]
What does \zh{水} mean? (\textquotedblleft{}Water\textquotedblright{}.) Say it once more.
\end{tcolorbox}
\section[\texorpdfstring{\zh{水} (shuǐ)}{水 (shuǐ)}]{\zh{水} — the middle first, then the sides}

\label{lesson:ZH-W39-shui}



\begin{quote}
The word first, then the character inside it.

\begin{itemize}
\raggedright
  \item \emph{Recall:} say \emph{shuǐ} — \textbf{R1}, one lesson back, before it decays
\end{itemize}

\begin{quote}
\zh{小} \zh{喝}
\end{quote}

\emph{Read it:} each one aloud, with nothing said first — \emph{xiǎo}, \emph{hē}
\end{quote}

\subsection*{Script — the middle first, then the sides}
\hlblockfigure{\detokenize{figures/ZH-W39-shui-filmstrip.pdf}}{How \zh{水} is written, stroke by stroke}

\begin{quote}
\zh{水}
\end{quote}

\textbf{One piece: the middle, then left, then right.}

\textbf{The middle, one stroke.} A long vertical that hooks to the left at the foot, as in \zh{小}.

\textbf{The left side, one stroke.} A short horizontal that turns and falls away to the lower left.

\textbf{The right side, two strokes.} A short left-falling stroke toward the middle, then a long right-falling stroke out to the lower right.

1 + 1 + 2 = four.

\subsection*{Your turn}
\begin{itemize}
\raggedright
  \item \emph{Write it:} \zh{水}, all four, with the model in front of you
  \item \emph{Cover:} write it again from memory, and count the pieces, not the strokes
\end{itemize}

\begin{tcolorbox}[breakable,colback=teal!4,colframe=teal!35!black,title={Before you move on}]
Which stroke of \zh{水} comes first? (\textbf{The hooked vertical down the middle}.) Which character you know also starts in the middle? (\textbf{\zh{小}}.) How many strokes on the right? (\textbf{Two}.)

Source: \href{https://raw.githubusercontent.com/chanind/hanzi-writer-data/68d10a4b21150cae5e1ebbd223eed289cf32d90c/data/\%E6\%B0\%B4.json}{Hanzi Writer Data \zh{水}, pinned stroke paths and medians}.
\end{tcolorbox}
\section[kě]{kě — thirsty}

\label{lesson:ZH-C39-ke}



\begin{quote}
Before the new one: say the Mandarin for to drink, then the Mandarin for water.
\end{quote}

\subsection*{You'll want to know: kě}
\emph{kě} — \textquotedblleft{}thirsty\textquotedblright{}.

\begin{quote}
\emph{nǐ kě ma} — Are you thirsty?
\end{quote}

Its characters come later; for now it is a word for the ear.

\begin{cousinweb}
\emph{kě} is inherited from Middle Chinese \emph{khat} (Guangyun, in Baxter's transcription), glossed 'hunger and thirst'.
\end{cousinweb}

\begin{grammarlens}[title={What you've built}]
One more word for everyday answers.
\end{grammarlens}

\subsection*{Your turn}
\begin{itemize}
\raggedright
  \item \emph{Say it:} \emph{kě}
  \item \emph{Say it:} \emph{kě}, once more
  \item \emph{Say it:} \emph{nǐ kě ma} — Are you thirsty?
  \item \emph{Recall:} say the Mandarin for to drink, then the Mandarin for water, then say \emph{kě} again
\end{itemize}

\begin{tcolorbox}[breakable,colback=teal!4,colframe=teal!35!black,title={Before you move on}]
What does kě mean? (\textquotedblleft{}Thirsty\textquotedblright{}.) Say it once more.
\end{tcolorbox}
\section[niúnǎi]{niúnǎi — milk}

\label{lesson:ZH-C39-niunai}



\begin{quote}
Before the new one: say the Mandarin for water, then the Mandarin for thirsty.
\end{quote}

\subsection*{You'll want to know: niúnǎi}
\emph{niúnǎi} — \textquotedblleft{}milk\textquotedblright{}.

\begin{quote}
\emph{háizi xǐhuan hē niúnǎi} — Children like milk.
\end{quote}

Its characters come later; for now it is a word for the ear.

\begin{cousinweb}
\emph{niúnǎi} is built from niú 'cattle' and nǎi 'milk'. The first is inherited from Middle Chinese \emph{ngjuw} (Guangyun, in Baxter's transcription), glossed 'a large beast'. The second is probably the word the Guangyun records as \emph{nreaX}, glossed 'milk', under another way of writing the character.
\end{cousinweb}

\begin{grammarlens}[title={What you've built}]
One more word for everyday answers.
\end{grammarlens}

\subsection*{Your turn}
\begin{itemize}
\raggedright
  \item \emph{Say it:} \emph{niúnǎi}
  \item \emph{Say it:} \emph{niúnǎi}, once more
  \item \emph{Say it:} \emph{háizi xǐhuan hē niúnǎi} — Children like milk.
  \item \emph{Recall:} say the Mandarin for water, then the Mandarin for thirsty, then say \emph{niúnǎi} again
\end{itemize}

\begin{tcolorbox}[breakable,colback=teal!4,colframe=teal!35!black,title={Before you move on}]
What does niúnǎi mean? (\textquotedblleft{}Milk\textquotedblright{}.) Say it once more.
\end{tcolorbox}
\section[chá]{chá — tea}

\label{lesson:ZH-C39-cha}



\begin{quote}
Before the new one: say the Mandarin for thirsty, then the Mandarin for milk.
\end{quote}

\subsection*{You'll want to know: chá}
\emph{chá} — \textquotedblleft{}tea\textquotedblright{}.

\begin{quote}
\emph{nǐ hē chá ma} — Would you like some tea?
\end{quote}

Its characters come later; for now it is a word for the ear.

\begin{cousinweb}
\emph{chá} is inherited from Middle Chinese \emph{drae} (Guangyun, in Baxter's transcription), which the dictionary explains as leaves gathered in spring and made into a drink.
\end{cousinweb}

\begin{grammarlens}[title={What you've built}]
One more word for everyday answers.
\end{grammarlens}

\subsection*{Your turn}
\begin{itemize}
\raggedright
  \item \emph{Say it:} \emph{chá}
  \item \emph{Say it:} \emph{chá}, once more
  \item \emph{Say it:} \emph{nǐ hē chá ma} — Would you like some tea?
  \item \emph{Recall:} say the Mandarin for thirsty, then the Mandarin for milk, then say \emph{chá} again
\end{itemize}

\begin{tcolorbox}[breakable,colback=teal!4,colframe=teal!35!black,title={Before you move on}]
What does chá mean? (\textquotedblleft{}Tea\textquotedblright{}.) Say it once more.
\end{tcolorbox}
